\section*{अध्याय \ref{ch:b2:powerseries} --- घात श्रेणियाँ}

\begin{solution}{exo:b2:powerseries:1}
$\sum \frac{n^2}{2^n}z^n$: अनुपात $\frac{(n+1)^2}{2^{n+1}}\cdot
\frac{2^n}{n^2} \to \frac12$: अतः $R = 2$।

$\sum \frac{z^n}{\binom{2n}{n}}$: $\binom{2n}{n} \sim
\frac{4^n}{\sqrt{\pi n}}$ (\cref{ex:b2:comparison:centralbinomial}), अतः बहुपद गुणकों को छोड़कर $\abs{a_n}^{-1} \approx 4^n$: $R = 4$ (अनुपात जाँच: $\frac{\binom{2n}{n}}{\binom{2n+2}{n+1}} =
\frac{(n+1)^2}{(2n+1)(2n+2)} \to \frac14$)।

$\sum z^{n!}$: गुणांक $k = n!$ होने पर $a_k = 1$, अन्यथा $0$। $\abs z < 1$ के लिए $\sum \abs z^{n!}$ अभिसरित होती है
(गुणोत्तर से प्रभुत्व-युक्त); और $\abs z \geq 1$ के लिए पद $0$ की ओर नहीं जाते: $R = 1$।

$\sum (2 + (-1)^n)^n z^n$: गुणांक $3^n$ (सम $n$) और $1$ (विषम $n$)। $a_nr^n$ की परिबद्धता के लिए $3r \leq 1$
आवश्यक है; और $r <
\frac13$ काम कर जाता है: $R = \frac13$।
\end{solution}

\begin{solution}{exo:b2:powerseries:2}
तीनों की त्रिज्या $1$ है (अनुपात जाँच)। $z = 1$ पर: $\sum 1$ अपसरित; $\sum\frac1n$ अपसरित;
$\sum\frac{1}{n^2}$ अभिसरित। और $z = -1$ पर: $\sum(-1)^n$ अपसरित; $\sum\frac{(-1)^n}{n}$ अभिसरित (एकांतर); $\sum\frac{(-1)^n}{n^2}$ अभिसरित
(\omterm{def:b2:series:def}{निरपेक्ष रूप से})। सीमा-व्यवहार त्रिज्या को दिखाई ही नहीं देता।
\end{solution}

\begin{solution}{exo:b2:powerseries:3}
$\frac{1}{1-x} = \sum x^n$ से, अवकलन कीजिए और $x$ से गुणा कीजिए (\cref{thm:b2:powerseries:calculus}):
\[
\sum n x^n = \frac{x}{(1-x)^2} .
\]
एक बार और अवकलन कीजिए और फिर $x$ से गुणा:
\[
\sum n^2 x^n = x\,\frac{\dd}{\dd x}\Bigl(\frac{x}{(1-x)^2}\Bigr)
= \frac{x(1 + x)}{(1-x)^3} .
\]
तीसरा योग: वह $-\ln(1 - x)$ का विषम भाग है:
\[
\sum_{n\geq0} \frac{x^{2n+1}}{2n+1}
= \frac{-\ln(1-x) + \ln(1+x)}{2}
= \frac12 \ln\frac{1+x}{1-x}
= \operatorname{artanh} x .
\]
\end{solution}

\begin{solution}{exo:b2:powerseries:4}
आंशिक भिन्न: $\frac{1}{(1-x)(2-x)} = \frac{1}{1-x} -
\frac{1}{2 - x} = \sum x^n - \frac12\sum \bigl(\frac x2\bigr)^n$:
\[
\frac{1}{(1-x)(2-x)} = \sum_{n\geq0} \Bigl(1 -
\frac{1}{2^{n+1}}\Bigr)x^n,
\qquad R = 1 .
\]

$\ln(1 + x + x^2) = \ln\frac{1 - x^3}{1 - x} = \ln(1 - x^3) -
\ln(1 - x) = \sum_{n\geq1}\frac{x^n}{n} -
\sum_{m\geq1}\frac{x^{3m}}{m}$: $x^n$ का गुणांक $3 \nmid n$ होने पर $\frac1n$ है, और $3 \mid n$ होने पर $\frac1n - \frac{3}{n} = -\frac2n$। त्रिज्या $1$
(निकटतम बाधा: $\ln(1-x^3)$ की श्रेणी)।
\end{solution}

\begin{solution}{exo:b2:powerseries:5}
$\sum_{m \geq 0} x^m$ (गुणांक $1$) और $\sum_{k\geq1} \frac{x^k}{k}$ (गुणांक $\frac1k$, $k \geq
1$) का \omterm{thm:b2:series:fubini}{कोशी गुणनफल}, जो दोनों $\abs x < 1$
के लिए \omterm{def:b2:series:def}{निरपेक्ष रूप से} अभिसारी हैं: गुणनफल में $x^n$ का गुणांक $\sum_{k=1}^{n} \frac1k \cdot 1 = H_n$ है। अतः
\[
\Bigl(\sum x^m\Bigr)\Bigl(\sum \frac{x^k}{k}\Bigr)
= \frac{1}{1-x}\cdot\bigl(-\ln(1-x)\bigr)
= \sum_{n\geq1} H_n x^n .
\]
\end{solution}

\begin{solution}{exo:b2:powerseries:6}
पुनरावृत्ति को $x^{n+1}$ से गुणा करके योग लीजिए ($\abs x <
\frac12$):
\[
U(x) - 1 = 2x\,U(x) + \sum_{n\geq0} n\,x^{n+1}
= 2x\,U(x) + \frac{x^2}{(1-x)^2} ,
\]
जिसमें \cref{exo:b2:powerseries:3} का उपयोग हुआ। अतः
\[
U(x) = \frac{1}{1 - 2x}\Bigl(1 + \frac{x^2}{(1-x)^2}\Bigr)
= \frac{1 - 2x + 2x^2}{(1-2x)(1-x)^2} .
\]
आंशिक भिन्न ($x = \frac12$ पर ढककर देखने से गुणांक $2$ मिलता है; दोहरे ध्रुव $x = 1$
पर गुणांक $-1$; और बीच वाला गुणांक $x = 0$ पर मूल्यांकन करने से लुप्त हो जाता
है):
\[
U(x) = \frac{2}{1-2x} - \frac{1}{(1 - x)^2} .
\]
दोनों का प्रसार करने पर:
\[
u_n = 2\cdot 2^n - (n + 1) = 2^{n+1} - n - 1 .
\]
(जाँच: $u_0 = 1$, $u_1 = 2u_0 + 0 = 2 = 4 - 2$।)
\end{solution}

\begin{solution}{exo:b2:powerseries:7}
आगमन: $f'(x) = \frac{2}{x^3}\eu^{-1/x^2}$, और यदि $f^{(n)}(x)
= P_n(\tfrac1x)\eu^{-1/x^2}$ तो
\[
f^{(n+1)}(x) =
\Bigl(-\frac{1}{x^2}\,P_n'\Bigl(\frac1x\Bigr) +
\frac{2}{x^3}\,P_n\Bigl(\frac1x\Bigr)\Bigr)\eu^{-1/x^2} :
\]
जो फिर कथित रूप का है। $0$ पर: $h \to 0$ होने पर अंतर-भागफल $\frac{f^{(n)}(h)}{h} =
\frac1h P_n(\frac1h)\eu^{-1/h^2} \to 0$, क्योंकि किसी
भी बहुपद $Q$ के लिए $u \to \pm\infty$ होने पर $Q(u)\,\eu^{-u^2} \to 0$ (घातांकी घातों को हरा देता है): अतः
आगमन से सारे $f^{(n)}(0)$ विद्यमान हैं और लुप्त हैं, और उसी सीमा से हर $f^{(n)}$ $0$
पर \omterm{def:b2:metric:continuity}{संतत} है। इस प्रकार $f \in C^\infty$, जिसकी $0$ पर टेलर श्रेणी शून्य है; और वह
टेलर श्रेणी $0 \neq f$ पर जुड़ती है: अतः $0$ पर \omterm{def:b2:powerseries:analytic}{वैश्लेषिक} नहीं।
\end{solution}

\begin{solution}{exo:b2:powerseries:8}
द्विपद श्रेणी: $\sqrt{1-4x} = \sum_{k\geq0}
\binom{1/2}{k}(-4x)^k$। $k = n + 1 \geq 1$ के लिए:
\begin{align*}
\binom{1/2}{n+1}(-4)^{n+1}
&= \frac{\frac12\bigl(\frac12 - 1\bigr)\cdots\bigl(\frac12 -
n\bigr)}{(n+1)!}\,(-4)^{n+1}\\
&= \frac{(-1)^n\,1\cdot3\cdots(2n-1)}{2^{n+1}(n+1)!}\,(-4)^{n+1}
= -\frac{2}{n+1}\cdot\frac{(2n)!}{n!\,n!} ,
\end{align*}
जिसमें $1\cdot3\cdots(2n-1) = \frac{(2n)!}{2^n n!}$ का उपयोग हुआ। अतः
\[
C(x) = \frac{1 - \sqrt{1-4x}}{2x}
= \frac{1}{2x}\sum_{n\geq0}\frac{2}{n+1}\binom{2n}{n}x^{n+1}
= \sum_{n\geq0} \frac{1}{n+1}\binom{2n}{n}\,x^n :
\]
$C_n = \frac{1}{n+1}\binom{2n}{n}$। त्रिज्या: $\frac14$ ($4x$ में द्विपद श्रेणी)। और \cref{ex:b2:comparison:centralbinomial} के द्वारा
अनंतस्पर्शिता:
\[
C_n \sim \frac{4^n}{\sqrt{\pi}\; n^{3/2}} .
\]
\end{solution}

\begin{solution}{exo:b2:powerseries:9}
$A_n = \sum_{k \leq n} a_k \to A$ के साथ: $0 \leq x < 1$ के लिए \omterm{thm:b2:series:abel}{आबेल योग} देता है
\[
\sum_{n=0}^{\infty} a_n x^n = (1 - x)\sum_{n=0}^{\infty} A_n x^n
\]
(दोनों पक्ष अभिसरित होते हैं: $(A_n)$ परिबद्ध है; और सर्वसमिका $a_n = A_n - A_{n-1}$ तथा पुनः
सूचीबद्ध करने से निकलती है)। चूँकि $(1 -
x)\sum x^n = 1$:
\[
\sum_n a_nx^n - A = (1-x)\sum_{n} (A_n - A)x^n .
\]
$\varepsilon$ दिया हो तो ऐसा $N$ चुनिए कि $n > N$ के लिए $\abs{A_n - A} \leq
\varepsilon$; तब
\[
\Bigl|\sum a_nx^n - A\Bigr|
\leq (1-x)\sum_{n \leq N}\abs{A_n - A} + \varepsilon(1 -
x)\sum_{n > N}x^n
\leq (1-x)\,C_N + \varepsilon ,
\]
और $x \to 1^-$ लेने पर: सभी $\varepsilon$ के लिए उच्चतम-सीमा $\leq \varepsilon$। अतः त्रिज्यीय सीमा
$A$ है।

अनुप्रयोग: $\sum \frac{(-1)^{n-1}}{n}$ अभिसरित होती है (एकांतर), और $x < 1$ के लिए उसकी घात श्रेणी
$\ln(1 + x)$ पर जुड़ती है: अतः आबेल से योग $\ln 2$ है। इसी प्रकार $\sum\frac{(-1)^n}{2n+1}x^{2n+1} =
\arctan x$ $x = 1$ पर $\frac\pi4$ देता
है --- अर्थात् प्रथम वर्ष की समाकल-उपपत्तियाँ, अब संरचनात्मक रूप में।
\end{solution}

\begin{solution}{exo:b2:powerseries:10}
$\sum\frac{x^n}{n}$ और $\sum\frac{x^n}{n+1}$ दोनों की त्रिज्या $1$ है, और $\frac{1}{n(n+1)} = \frac1n - \frac1{n+1}$, अतः $0
< \abs x < 1$ के लिए:
\[
\sum_{n\geq1}\frac{x^n}{n(n+1)}
= -\ln(1-x) - \frac1x\sum_{n\geq1}\frac{x^{n+1}}{n+1}
= -\ln(1-x) - \frac{-\ln(1-x) - x}{x}
= 1 + \frac{1-x}{x}\ln(1-x) .
\]
त्रिज्या $1$; और $\norm{x^n/(n(n+1))}_{\infty,\intcc{-1}1} =
\frac{1}{n(n+1)}$ \omterm{def:b2:series:summable}{योग्य} है: अतः $\intcc{-1}{1}$ पर \omterm{def:b2:funcseq:series}{सामान्य अभिसरण}, इसलिए योग वहाँ
\omterm{def:b2:metric:continuity}{संतत} है। $x \to
1^-$ होने पर $(1-x)\ln(1-x) \to 0$ और बंद रूप $1$ की ओर जाता है --- जो $x = 1$ पर
क्रमिक निरसन मान $\sum\frac{1}{n(n+1)} =
\lim_N\bigl(1 - \frac{1}{N+1}\bigr) = 1$ के अनुरूप है।
\end{solution}

\begin{solution}{exo:b2:powerseries:11}
$n!$ क्रमचयों को उनके अचर-बिंदु समुच्चय के अनुसार छाँटने पर: $k$ अचर बिंदु
चुनना ($\binom nk$ प्रकार से) और शेष $n - k$ वस्तुओं का अव्यवस्थापन करना $n! = \sum_{k=0}^n\binom nk D_{n-k}$ देता है।
$n!$ से भाग देने पर:
\[
1 = \sum_{k=0}^{n}\frac{1}{k!}\cdot\frac{D_{n-k}}{(n-k)!} ,
\]
जो ठीक यह कहता है कि $\eu^x =
\sum\frac{x^k}{k!}$ और $D(x) = \sum D_n\frac{x^n}{n!}$ का \omterm{thm:b2:series:fubini}{कोशी गुणनफल} $\sum x^n = \frac{1}{1-x}$ है। दोनों गुणनखंड $\abs x < 1$
के लिए \omterm{def:b2:series:def}{निरपेक्ष रूप से} अभिसरित होते हैं ($D_n \leq n!$, अतः $D$ गुणोत्तर श्रेणी से
प्रभुत्व-युक्त है): इसलिए गुणनफल सर्वसमिका वैध है (\cref{prop:b2:powerseries:operations}), और
\[
D(x) = \frac{\eu^{-x}}{1-x} .
\]
$\eu^{-x} = \sum\frac{(-1)^kx^k}{k!}$ तथा $\sum x^m$ का \omterm{thm:b2:series:fubini}{कोशी गुणनफल}: $x^n$ का गुणांक $\sum_{k=0}^{n}\frac{(-1)^k}{k!}$ है, और घात श्रेणी के
गुणांकों की अद्वितीयता (\cref{thm:b2:powerseries:calculus}) से:
\[
\frac{D_n}{n!} = \sum_{k=0}^{n}\frac{(-1)^k}{k!}
\xrightarrow[n\to\infty]{} \eu^{-1} :
\]
अर्थात् $n$ चाहे जो हो, लगभग $37\%$ क्रमचय अव्यवस्थापन होते हैं।
\end{solution}

\begin{solution}{exo:b2:powerseries:12}
$-4x$ पर $\alpha = -\frac12$ वाली द्विपद श्रेणी:
\[
\binom{-1/2}{n}(-4)^n
= \frac{\bigl(-\frac12\bigr)\bigl(-\frac32\bigr)\cdots
\bigl(-\frac{2n-1}2\bigr)}{n!}(-4)^n
= \frac{1\cdot3\cdots(2n-1)}{2^n\,n!}\,4^n
= \frac{(2n)!}{2^n n!}\cdot\frac{2^n}{n!}
= \binom{2n}{n},
\]
जिसमें $1\cdot3\cdots(2n-1) = \frac{(2n)!}{2^nn!}$ का उपयोग हुआ। अतः $\abs{4x} < 1$ के लिए $(1-4x)^{-1/2} = \sum\binom{2n}nx^n$। वर्ग करने पर (\omterm{thm:b2:series:fubini}{कोशी गुणनफल},
जो निरपेक्ष अभिसरण से वैध है) और $\frac{1}{1-4x} = \sum 4^nx^n$ से तुलना करने पर: वर्ग में $x^n$ का
गुणांक $\sum_{k=0}^n\binom{2k}k\binom{2n-2k}{n-k}$ है, और गुणांकों की अद्वितीयता देती है
\[
\sum_{k=0}^{n}\binom{2k}{k}\binom{2n-2k}{n-k} = 4^n .
\]
\end{solution}

\begin{solution}{pb:b2:powerseries:1}
\textbf{1.} $a_n = r_n - r_{n+1}$ के साथ, खंडशः योग:
\[
\sum_{n=N}^{M} a_nx^n
= r_Nx^N + \sum_{n=N+1}^{M} r_n\bigl(x^n - x^{n-1}\bigr)
- r_{M+1}x^M .
\]
$0 \leq x \leq 1$ के लिए वृद्धियाँ $x^{n-1} - x^n$ अऋणात्मक हैं और क्रमिक निरसन से $x^N - x^M$ तक जुड़ जाती
हैं; और $s = \sup_{n\geq
N}\abs{r_n}$ के साथ:
\[
\Bigl|\sum_{n=N}^{M}a_nx^n\Bigr|
\leq s\bigl(x^N + (x^N - x^M) + x^M\bigr) = 2s\,x^N \leq 2s .
\]

\textbf{2.} चूँकि $r_n \to 0$, $\sup_{n\geq N}\abs{r_n} \to 0$: अतः प्रश्न 1 ठीक $\intcc{0}{1}$ पर एकसमान कोशी कसौटी है,
इसलिए $\sum a_nx^n$ वहाँ \omterm{def:b2:funcseq:def}{एकसमान रूप से} अभिसरित होती है और उसका योग \omterm{def:b2:metric:continuity}{संतत} है (\cref{thm:b2:funcseq:seriestransfer})।
और $1$ पर मान $\sum a_n$ होने से $1$ पर संततता \cref{exo:b2:powerseries:9} की त्रिज्यीय सीमा ही है।

\textbf{3.} $\abs x < 1$ के लिए तीनों घात श्रेणियाँ \omterm{def:b2:series:def}{निरपेक्ष रूप से} अभिसरित होती
हैं और $\bigl(\sum a_nx^n\bigr)\bigl(\sum b_nx^n\bigr)
= \sum c_nx^n$ (\cref{prop:b2:powerseries:operations})। प्रश्न 2 से हर गुणनखंड तथा गुणनफल पक्ष $\intcc{0}{1}$ पर \omterm{def:b2:metric:continuity}{संतत} हैं
(उनकी गुणांक-श्रेणियाँ परिकल्पना से अभिसरित होती हैं); और सर्वसमिका में
$x \to 1^-$ लेने पर: $AB = C$।

\textbf{4.} यहाँ समांतर--गुणोत्तर माध्य से
\[
\abs{c_n} = \sum_{k=0}^{n}
\frac{1}{\sqrt{(k+1)(n-k+1)}}
\geq \sum_{k=0}^{n}\frac{2}{n+2}
= \frac{2(n+1)}{n+2} \geq 1,
\]
अर्थात् $\sqrt{(k+1)(n-k+1)} \leq \frac{(k+1) +
(n-k+1)}{2} = \frac{n+2}{2}$। $\sum c_n$ का व्यापक पद $0$ की ओर नहीं जाता: अतः \omterm{thm:b2:series:fubini}{कोशी गुणनफल}
अपसरित होता है, यद्यपि दोनों गुणनखंड अभिसरित होते हैं (एकांतर श्रेणियाँ)।

\textbf{5.} \cref{exo:b2:powerseries:5} देता है $\frac{-\ln(1-x)}{1-x} = \sum H_nx^n$ ($\abs x < 1$)। पद-दर-पद प्रतिअवकलज (\cref{thm:b2:powerseries:calculus}~(2)), और
दोनों पक्ष $0$ पर लुप्त:
\[
\frac{\bigl(\ln(1-x)\bigr)^2}{2}
= \sum_{n\geq1}\frac{H_n}{n+1}\,x^{n+1} .
\]
ह्रास: $(n+2)H_n \geq (n+1)H_{n+1}$ का अर्थ $H_n \geq
1$ है, जो $n \geq 1$ के लिए सत्य है; और $\frac{H_n}{n+1} \sim \frac{\ln
n}{n} \to 0$: $x = -1$ पर श्रेणी
एकांतर जाँच से अभिसरित होती है। $x \mapsto -x$ प्रतिस्थापित करके $x = 1$ पर प्रश्न 2
लगाने पर:
\[
\frac{(\ln 2)^2}{2}
= \sum_{n\geq1}\frac{H_n}{n+1}(-1)^{n+1},
\]
अर्थात् घोषित मान।

\textbf{6.} $x \to 1^-$ होने पर $f(x) = \sum(-1)^nx^n = \frac{1}{1+x} \to \frac12$: अतः $\frac12$ पर आबेल-योग्य। परंतु आंशिक योग $1, 0, 1, 0, \dots$
हैं: अर्थात् अपसारी।

\textbf{7.} $\varepsilon > 0$ दिया हो तो ऐसा $m$ चुनिए कि $n > m$ के लिए $\abs{u_n}
\leq \varepsilon$; और $N \geq m$ के
लिए:
\[
\Bigl|\frac{u_1 + \dots + u_N}{N}\Bigr|
\leq \frac{\abs{u_1} + \dots + \abs{u_m}}{N} +
\varepsilon\,\frac{N - m}{N}
\leq \frac{C_m}{N} + \varepsilon,
\]
अतः प्रत्येक $\varepsilon$ के लिए $\limsup \leq \varepsilon$: इसलिए माध्य $0$ की ओर जाते हैं।

\textbf{8.} $0 \leq x \leq 1$ के लिए: $1 - x^n = (1-x)(1 + x +
\dots + x^{n-1}) \leq n(1-x)$, अतः
\[
\Bigl|\sum_{n=0}^N a_n(1 - x_N^n)\Bigr|
\leq (1 - x_N)\sum_{n=1}^N n\abs{a_n}
= \frac1N\sum_{n=1}^{N}n\abs{a_n} .
\]
और $n > N$ के लिए: $\abs{a_n} = \frac{n\abs{a_n}}{n} \leq
\frac{1}{N}\sup_{m>N}m\abs{a_m}$, तथा $\sum_{n>N}x_N^n \leq
\frac{1}{1 - x_N} = N$:
\[
\Bigl|\sum_{n>N}a_nx_N^n\Bigr|
\leq \frac{\sup_{m>N}m\abs{a_m}}{N}\cdot N
= \sup_{m>N}\,m\abs{a_m} .
\]

\textbf{9.} अपघटन कीजिए
\[
A_N - L = \sum_{n=0}^{N}a_n\bigl(1 - x_N^n\bigr)
- \sum_{n>N}a_nx_N^n + \bigl(f(x_N) - L\bigr) .
\]
पहला पद प्रश्न 7 से $0$ की ओर जाता है ($n\abs{a_n} \to 0$ के माध्य), दूसरा प्रश्न 8 से
(उच्चतम $0$ की ओर जाता है), और तीसरा इसलिए कि $x_N \to 1^-$ तथा $f(x) \to L$। अतः $A_N \to L$: यही
टाउबर की प्रमेय है।

\textbf{10.} $x \in \intco{0}{1}$ और किसी भी $N$ के लिए: $\sum_{n\leq N}a_nx^n \leq f(x) \leq M$ (अऋणात्मक पद)। परिमित योग में
$x \to 1^-$ लीजिए: $\sum_{n\leq N}a_n \leq M$। आंशिक योग वर्धमान और परिबद्ध हैं: अतः $\sum a_n$ अभिसरित होती है,
और फिर प्रश्न 2 $\lim_{x\to1^-}f(x) =
\sum a_n$ दे देता है।

\textbf{11.} $\sigma_N - L$ $N$ संख्याओं $A_n - L$ ($0 \leq n < N$) का माध्य है, और वे $0$ की ओर जाती
हैं: अतः प्रश्न 7।

\textbf{12.} सम $n$ के लिए $A_n = 1$, विषम के लिए $0$: $A_0 + \dots
+ A_{N-1} = \lceil N/2\rceil$, अतः $\sigma_N =
\frac{\lceil N/2\rceil}{N} \to \frac12$, जो प्रश्न
6 का आबेल मान है।

\textbf{13.} $\sigma_N \to L$ के अंतर्गत $S_n = O(n)$, अतः $A_n = S_n - S_{n-1} = O(n)$ और $a_n = O(n)$: इसलिए नीचे की सारी
श्रेणियों की त्रिज्या $\geq 1$ है। $\abs x < 1$ के लिए $a_n = A_n -
A_{n-1}$ और $A_nx^n \to 0$ से:
\[
(1-x)\sum_n A_nx^n = \sum_n A_nx^n - \sum_n A_nx^{n+1}
= \sum_n a_nx^n = f(x),
\]
और सर्वथा $(1-x)\sum S_nx^n = \sum A_nx^n$, अतः $f(x) =
(1-x)^2\sum_n S_nx^n = (1-x)^2\sum_n(n+1)\sigma_{n+1}x^n$। अंत में $\sum(n+1)x^n = \frac{1}{(1-x)^2}$ (\cref{exo:b2:powerseries:3}), जो दूसरी सर्वसमिका है।

\textbf{14.} पहली सर्वसमिका में से दूसरी का $L$ गुना घटाने पर:
\[
f(x) - L = (1-x)^2\sum_{n\geq0}(n+1)
\bigl(\sigma_{n+1} - L\bigr)x^n .
\]
$\varepsilon$ दिया हो तो ऐसा $N$ चुनिए कि $n \geq N$ के लिए $\abs{\sigma_{n+1} - L}
\leq \varepsilon$; तब
\[
\abs{f(x) - L} \leq (1-x)^2 C_N +
\varepsilon(1-x)^2\sum_{n}(n+1)x^n
= (1-x)^2C_N + \varepsilon ,
\]
और $x \to 1^-$ लेने पर: $\limsup \leq \varepsilon$। अतः $f(x) \to L$: यही फ्रोबेनियस की प्रमेय है।

\textbf{15.} $f(x) = \sum(-1)^n(n+1)x^n = \frac{1}{(1+x)^2}$ ($-x$ पर गुणोत्तर श्रेणी का अवकलन): अतः आबेल मान $\frac14$।
आंशिक योग: $A_{2k} = k+1$, $A_{2k+1} = -(k+1)$ (तत्काल आगमन)। तब $S_{2m-1} = 0$ (क्रमागत युग्म कट जाते हैं)
और $S_{2m} = m + 1$, अतः
\[
\sigma_{2m} = \frac{S_{2m-1}}{2m} = 0,
\qquad
\sigma_{2m+1} = \frac{m+1}{2m+1} \to \frac12 :
\]
$(\sigma_N)$ के दो भिन्न गुच्छन मान हैं: अर्थात् चेज़ारो-योग्य नहीं। प्रश्न 6 तथा
11--14 के साथ पदानुक्रम अभिसारी $\Rightarrow$ चेज़ारो $\Rightarrow$ आबेल दोनों तीरों पर
यथार्थ है।

\textbf{16.} प्रतिअवकलज श्रेणी $F(x) =
\sum\frac{a_n}{n+1}x^{n+1}$ की त्रिज्या वही है और $\intco{0}{1}$ पर $F' = f$
(\cref{thm:b2:powerseries:calculus}); और चूँकि $\sum\frac{a_n}{n+1}$ अभिसरित होती है, भाग I (प्रश्न 2) $F$ को $\intcc{0}{1}$ पर \omterm{def:b2:metric:continuity}{संतत}
बना देता है। और $\int_0^x f = F(x)$ होने से (अवकलज बराबर, $0$ पर मान $0$ बराबर),
\[
\int_0^x f \xrightarrow[x\to1^-]{} F(1)
= \sum_{n\geq0}\frac{a_n}{n+1} :
\]
अर्थात् \omterm{def:b2:integration:improper}{अनुचित समाकल} कथित मान के साथ विद्यमान है।

\textbf{17.} $\frac{\ln(1+x)}{x} =
\sum_{n\geq1}\frac{(-1)^{n-1}}{n}x^{n-1}$ (त्रिज्या $1$; $0$ पर \omterm{def:b2:metric:continuity}{संतत})। $\frac{a_m}{m+1}$ की श्रेणी $\sum_{n\geq1}\frac{(-1)^{n-1}}{n^2}$ है, जो
\omterm{def:b2:series:def}{निरपेक्ष रूप से} अभिसारी है: अतः प्रश्न 16 $\int_0^1\frac{\ln(1+x)}{x}\dd x = \eta$ देता है। और \omterm{def:b2:series:def}{निरपेक्ष रूप
से} अभिसारी $\sum\frac1{n^2}$ में सम तथा विषम पुनः समूहबद्ध कीजिए:
\[
\eta = \sum_{\text{विषम}}\frac1{n^2} -
\sum_{\text{सम}}\frac1{n^2}
= \sum_{n}\frac1{n^2} - 2\sum_{k}\frac1{(2k)^2}
= \Bigl(1 - \frac12\Bigr)\sum_n\frac1{n^2}
= \frac12\sum_{n\geq1}\frac1{n^2} .
\]

\textbf{18.} लाइब्नित्स: $\frac{1}{3n+1}\downarrow0$, अतः श्रेणी अभिसरित होती है। $\intco{0}{1}$ पर $\sum(-1)^nx^{3n} =
\frac{1}{1+x^3}$, और
$\sum\frac{(-1)^n}{3n+1}$ अभिसरित होती है: अतः प्रश्न 16 $\sum\frac{(-1)^n}{3n+1} =
\int_0^1\frac{\dd x}{1+x^3}$ देता है। आंशिक भिन्न (जाँच: $\frac13(x^2-x+1) + \frac{2-x}{3}(1+x) = 1$):
\[
\int_0^1\frac{\dd x}{1+x^3}
= \frac13\ln2 + \frac13\int_0^1\frac{2-x}{x^2-x+1}\dd x .
\]
$2 - x = -\frac12(2x-1) + \frac32$ लिखने पर: $\ln(x^2-x+1)$ वाला भाग दोनों सिरों पर लुप्त हो जाता है, और
\[
\frac32\int_0^1\frac{\dd x}{(x-\frac12)^2 + \frac34}
= \frac32\cdot\frac{2}{\sqrt3}
\Bigl[\arctan\frac{2x-1}{\sqrt3}\Bigr]_0^1
= \sqrt3\cdot\frac{\pi}{3} = \frac{\pi}{\sqrt3} .
\]
कुल: $\frac13\bigl(\ln2 + \frac{\pi}{\sqrt3}\bigr)$।

\textbf{19.} \cref{exo:b2:powerseries:12} की श्रेणी में $t = x^2$ प्रतिस्थापित करके पद-दर-पद समाकलन
कीजिए ($0$ पर लुप्त होने वाला $(1-x^2)^{-1/2}$ का प्रतिअवकलज $\arcsin$ है):
\[
\arcsin x = \sum_{n\geq0}
\frac{\binom{2n}n}{4^n(2n+1)}x^{2n+1}
\qquad(\abs x < 1) .
\]
गुणांक $\sim \frac{1}{2\sqrt\pi\,n^{3/2}}$ हैं (\cref{ex:b2:comparison:centralbinomial}), जो \omterm{def:b2:series:summable}{योग्य} हैं: अतः श्रेणी $\intcc{-1}{1}$ पर \emph{सामान्य} रूप
से अभिसरित होती है, उसका योग वहाँ \omterm{def:b2:metric:continuity}{संतत} है, और $\intoo{-1}{1}$ पर \omterm{def:b2:metric:continuity}{संतत} $\arcsin$ से मेल
खाता है, इसलिए $x = 1$ पर भी:
\[
\sum_{n\geq0}\frac{\binom{2n}n}{4^n(2n+1)}
= \arcsin 1 = \frac\pi2 .
\]

\textbf{20.} $C_n4^{-n} \sim \frac{1}{\sqrt\pi\,n^{3/2}}$ (\cref{exo:b2:powerseries:8}): $\intcc{0}{\frac14}$ पर $\sum
C_nx^n$ का \omterm{def:b2:funcseq:series}{सामान्य अभिसरण}, अतः उसका योग वहाँ
\omterm{def:b2:metric:continuity}{संतत} है; और $\intoo{0}{\frac14}$ पर वह
$\frac{1-\sqrt{1-4x}}{2x}$
के बराबर है (\cref{ex:b2:powerseries:catalan}), जिसकी $\frac14^-$ पर सीमा $\frac{1-0}{1/2} = 2$ है। अतः $\sum_{n\geq0}
C_n4^{-n} = 2$।

\textbf{21.} $g(x) = \sum_{k\geq0}a_{m+k}x^k$ के साथ $f(x) = \sum_{n\geq m}a_nx^n = x^m g(x)$; और यदि $(a_nr^n)$ परिबद्ध हो तो $(a_{m+k}r^k)$ भी ($r^m$ से भाग
दीजिए): अतः $g$ की त्रिज्या $\geq
R$ है। $g$ $g(0) = a_m \neq 0$ के साथ \omterm{def:b2:metric:continuity}{संतत} है, अतः किसी $\intcc{-\delta}{\delta}$
पर $g \neq 0$, और $0 < \abs x \leq \delta$ के लिए $f(x) = x^mg(x) \neq 0$।

\textbf{22.} $d = f - h$ $0$ के पास किसी घात श्रेणी का योग है जो अशून्य बिंदुओं
$x_k \to 0$ पर लुप्त होता है। यदि $d$ का कोई गुणांक अशून्य होता, तो प्रश्न 21 $0$
का ऐसा छिद्रित प्रतिवेश देता जिसमें $d$ का कोई शून्य न होता --- जो $d(x_k) = 0$ के
विरुद्ध है। अतः $d$ के सारे गुणांक लुप्त हैं: इसलिए $f$ और $h$ के गुणांक
बराबर हैं और वे $0$ के पास मेल खाते हैं।

\textbf{23.} फलन $h(x) = \frac{1}{1+x^2} =
\sum(-1)^nx^{2n}$ (त्रिज्या $1$) $h(\frac1k) =
\frac{1}{1 + 1/k^2} = \frac{k^2}{k^2+1}$ पूरा करता है। समान मानों वाला कोई
भी \omterm{def:b2:powerseries:analytic}{वैश्लेषिक} $f$ बिंदुओं $\frac1k \to 0$ पर $h$ से मेल खाता है: अतः तादात्म्य प्रमेय
(प्रश्न 22) से $0$ के पास $f = \frac{1}{1+x^2}$ --- अर्थात् अद्वितीय हल।

\textbf{24.} मान लीजिए $Z$ $I$ के उन बिंदुओं का समुच्चय है जिनका कोई ऐसा
प्रतिवेश हो जिस पर $f$ सर्वथा लुप्त हो: यह परिभाषा से \omterm{def:b2:metric:topology}{विवृत} और अरिक्त है
(वह उपअंतराल)। $I$ में संवृत भी: यदि $y
\in I$ $Z$ के बिंदुओं की सीमा हो, तो $y$
$f$ के शून्यों का संचय-बिंदु है; और $y$ पर $f$ का घात श्रेणी में प्रसार
करके ($y$ पर पुनःकेंद्रित) प्रश्न 21--22 लगाने पर $y$ पर सारे गुणांक
लुप्त हो जाते हैं, अतः $y$ के पास $f \equiv 0$: $y
\in Z$। अंतराल \omterm{def:b2:metric:connected}{संबद्ध} है, अतः $Z = I$:
$I$ पर $f \equiv 0$। विशेष रूप से $\R$ पर \omterm{def:b2:powerseries:analytic}{वैश्लेषिक फलन} जो किसी \omterm{def:b2:metric:compact}{संहत} समुच्चय के
बाहर लुप्त हो, किसी अंतराल पर लुप्त होता है, अतः सर्वत्र: यानी कोई अशून्य
\omterm{def:b2:powerseries:analytic}{वैश्लेषिक} उभार नहीं। $C^\infty$ की दुनिया भिन्न है: \cref{exo:b2:powerseries:7} के चपटे फलन को जोड़कर
(जैसे $x \mapsto
\eu^{-1/x^2}\mathbf 1_{x>0}$ और उसका दर्पण) \omterm{def:b2:metric:compact}{संहत} वाहक वाले चिकने उभार बन जाते हैं।

\textbf{25.} (क) \omterm{def:b2:funcseq:series}{सामान्य अभिसरण} चक्रिका के यथार्थतः भीतर की \omterm{def:b2:metric:compact}{संहत}
उपचक्रिकाओं पर बसता है; और आबेल की प्रमेय संततता को किसी सीमा-बिंदु तक
बढ़ा देती है, केवल इस परिकल्पना पर कि गुणांक-श्रेणी वहाँ अभिसरित हो। (ख)
विलोम टाउबर की शर्त $na_n \to 0$ के अंतर्गत सत्य है (प्रश्न 9), और चेज़ारो योग्यता
अभिसरण तथा आबेल योग्यता के यथार्थतः बीच में बैठती है (फ्रोबेनियस, प्रश्न
14--15)। (ग) किसी मित्र के लिए: गुणोत्तर श्रेणी को सीमा तक समाकलित करके
और फिर आंशिक भिन्न लेकर $\sum\frac{(-1)^n}{3n+1} = \int_0^1\frac{\dd x}{1+x^3}$। (घ) चेज़ारो माध्य फूरिये अध्याय में फेयेर की
प्रमेय के रूप में लौटते हैं, जहाँ आंशिक योगों का औसत लेने से \omterm{def:b2:funcseq:def}{बिंदुवार
अभिसरण} की विफलता ठीक हो जाती है --- वही दवा, नया रोगी।
\end{solution}
